\section{Introduction}
\label{sec:introduction}

Plant functional traits are measurable properties that influence how organisms respond to and affect their environment, and are central to understanding ecosystem functioning, biodiversity, and the global carbon cycle \citep{kattge2020try}. \jl{In particular,} traits such as leaf chlorophyll content, carotenoid concentration, leaf mass per area, and leaf area index govern photosynthetic capacity, light interception, and nutrient cycling, making their accurate quantification essential for ecological research and precision agriculture \citep{homolova2013review}. \jl{To meet this need,} hyperspectral reflectance spectroscopy offers a non-destructive pathway for estimating these traits at scale, since foliar biochemistry and canopy structure imprint characteristic absorption features across the 400--2500\,nm spectral range \citep{curran1989imaging, fu2020photosynthetic}. \sm{SUENA MUY CORTADO COMO FRASE TRAS FARASE SIN CONEXION}

Traditional approaches for deriving traits from spectra rely on vegetation indices, look-up table inversion of radiative transfer models such as PROSAIL \citep{feret2021prospectpro, ferreira2025enhancing}, or statistical methods---principally partial least squares regression (PLSR)---that project the high-dimensional spectral space onto a low-dimensional latent representation \citep{fu2020photosynthetic, ji2024unveiling}. While effective, PLSR models are typically trained and validated on individual traits, require careful feature engineering, and can struggle to capture the non-linear relationships between reflectance and biochemical concentrations that arise from multiple scattering and overlapping absorption features \citep{verhoef2007sail}. More recently, one-dimensional convolutional neural networks (1D-CNNs) have been applied directly to spectral vectors, learning hierarchical features end-to-end and achieving state-of-the-art multi-trait retrieval performance without manual feature selection \citep{cherif2023spectra, furbank2021wheat}. \citet{cherif2023spectra} demonstrated that a single EfficientNet-based 1D-CNN trained simultaneously on 20 traits across 50 heterogeneous field datasets outperformed PLSR baselines, establishing a benchmark for transferable multi-trait models. Their subsequent \texttt{GreenHyperSpectra} dataset \citep{cherif2025greenhyperspectra} further advanced this line of work by incorporating semi-supervised and self-supervised pretraining strategies, with a 1D masked autoencoder (MAE) achieving the best results ($R^2 = 0.645$) after fine-tuning on labeled data.

Despite this progress, 1D-CNNs process the spectrum as a flat sequence, inherently limiting their receptive field to local neighbourhoods of adjacent bands. Spectral signatures of plant traits, however, arise from the interplay of multiple absorption features distributed across the full wavelength range---for example, the joint influence of pigment absorption in the visible, water absorption in the shortwave infrared, and structural scattering near the red edge \citep{curran1989imaging, homolova2013review}. Capturing these long-range inter-band dependencies within a 1D framework requires either very deep networks or large kernel sizes, both of which increase the risk of overfitting when labeled training data are scarce.

An alternative strategy, increasingly explored in signal processing and chemometrics, is to transform one-dimensional signals into two-dimensional image representations and leverage the powerful feature extraction capabilities of 2D-CNNs originally developed for computer vision \citep{wang2015gaf, garcia2022temporal}. Several encoding methods have been proposed for this purpose. Gramian Angular Fields (GAF) and Markov Transition Fields (MTF) encode temporal or spectral correlations as square images by mapping signal values to angular or transition-probability spaces \citep{wang2015gaf, ganesan2021gaf}. Continuous wavelet transform (CWT) scalograms decompose signals across multiple frequency scales, producing time-frequency representations that have proven effective for fault diagnosis, medical signal classification, and spectral analysis \citep{addison2017wavelet, Lopatin2023-fn, ong2025sugarcane, zhuang2023cwt}. Short-time Fourier transform spectrograms offer a complementary multi-window decomposition \citep{shuai2025multichannel}, while direct reshaping---simply folding the 1D vector into a 2D matrix---has been shown to produce surprisingly competitive results by preserving wavelength adjacency as spatial texture \citep{hennessy2022reshaping, deev2024spectrum, sharma2019deepinsight}. More recently, these encoding techniques have begun to appear in spectroscopic applications, including GAF-based adulteration detection in food products \citep{zhang2025gaf} and CWT-assisted disease recognition from near-infrared spectra \citep{ong2025sugarcane}.

However, the application of 1D-to-2D spectral transformations to vegetation trait retrieval from hyperspectral data remains largely unexplored. Existing studies in this domain have predominantly focused on classification tasks (e.g., species discrimination) rather than continuous regression of biophysical variables \citep{hennessy2022reshaping}, and no systematic comparison of multiple transformation methods under controlled experimental conditions has been conducted. Furthermore, the conversion of spectra to 2D images unlocks capabilities exclusive to the vision domain that are inaccessible to 1D spectral architectures. \jl{The most important of these is} the ability to leverage self-supervised pretraining with 2D masked autoencoders \citep{he2022mae}, which have demonstrated remarkable success in learning visual representations from unlabeled image data. While 1D MAEs have been applied to spectral pretraining \citep{cherif2025greenhyperspectra}, the potential of 2D MAEs operating on transformed spectral images has not been investigated.

In this study, we aim to: (i) systematically compare the performance of six 1D-to-2D spectral transformation methods as spectral encoding strategies for 2D-CNN-based trait retrieval; and (ii) assess whether converting spectra to 2D images enables self-supervised pretraining via masked autoencoders to improve trait predictions beyond what 1D pretraining achieves. We used the \texttt{GreenHyperSpectra} dataset to benchmark our 2D approach against the state-of-the-art 1D methods \citep{cherif2025greenhyperspectra}. 
